% ═════════════════════════════════════════════════════════════════════════════
\section{Cutoff Methods and Dynamic Cutoffs}
\label{sec:cutoff}
% ═════════════════════════════════════════════════════════════════════════════

A central idea in parameterized verification is the \emph{cutoff theorem}.
The observation behind it is that, for certain classes of systems, one does
not need to verify infinitely many instances. Instead, it
suffices to verify systems up to some finite bound, the cutoff, and
correctness for all larger instances follows automatically~\cite{EN95,EK00}.

Intuitively, this works because processes in many symmetric systems are
interchangeable. If the behavior of a system with~$n$ processes can be
faithfully represented by a system with at most~$k$ processes (the cutoff),
then verifying the $k$-process system is enough. Cutoff results have been
established for important classes such as token-passing rings, guarded
protocols, and systems where processes communicate in restricted ways. This
section first presents such \emph{static} cutoffs, and then \emph{dynamic}
cutoffs, which are discovered during verification.

\subsection{Static Cutoffs}

We say that a number~$k$ is a \emph{cutoff} for a class of parameterized
systems and properties if, for every system and every property in the class,
the property holds for all instances of the system whenever it holds for all
instances of size at most~$k$. A cutoff theorem reduces parameterized
verification to a finite number of finite-state model checking problems, which
can be solved with standard tools. The cutoff typically depends on the class
of systems, on the number of process indices mentioned in the property, and
sometimes on the size of the process template.

\paragraph{Token rings.}
Emerson and Namjoshi~\cite{EN95} studied rings of identical processes in which
a single token, carrying no data, is passed from each process to its neighbor.
They showed that properties in indexed temporal logic that refer to a fixed
number of processes have a cutoff that is a small constant, independent of the
process template. The reason is that a process in a token ring observes the
rest of the ring only through the arrival and departure of the token, and from
the point of view of a few processes, a large ring behaves like a small one in
which the remaining processes have been merged. The result was later extended
from rings to other network topologies by decomposing a network into small
fragments~\cite{CTTV04}, and the existence of cutoffs for token-passing systems
has been characterized in terms of the topology and the specification
logic~\cite{AJKR14}.

\paragraph{Guarded protocols.}
Emerson and Kahlon~\cite{EK00} considered systems of identical processes whose
transitions are guarded by conditions on the other processes. A guard is
\emph{conjunctive} if it requires all other processes to be in a given set of
local states, and \emph{disjunctive} if it requires at least one other process
to be in such a set. They established cutoffs that grow linearly with the
number of local states of the template. For disjunctive guards, the argument
is a simulation: in a run of a large instance, each guard needs only one
witness, and for each local state that occurs in the run it suffices to keep
one process that reaches this state and then stays there, so that it can serve
as a witness whenever needed. A small instance with these extra processes can
reproduce the behavior of the processes mentioned in the property. For conjunctive guards, the argument uses projection instead:
removing processes can only make a conjunctive guard easier to satisfy, so the
projection of a run of a large instance onto a subset of its processes is a
run of a smaller instance. For safety properties this argument is direct,
while liveness properties require more care, since the projection of a fair
run need not be fair, and deadlocks of the large instance need not appear in
the small one. The cutoffs were later sharpened, and extended to fairness and
deadlock detection~\cite{AJK16}, and specialized cutoff methods were developed
for cache coherence protocols~\cite{EK03}.

The running example of Section~\ref{sec:running} is a guarded protocol: the
transition from~$\Wait$ to~$\Crit$ has the conjunctive guard that all other
processes are in $\Idle$ or~$\Wait$. If two processes are in~$\Crit$ at the
same time in a run of some instance, then the projection of the run onto these
two processes is a run of the instance of size two that violates mutual
exclusion. Checking the instance of size two, which has only nine
configurations, is therefore enough to establish mutual exclusion for all
sizes.

\paragraph{Limitations.}
Static cutoff theorems are attractive because they reduce the problem to
standard model checking. However, the bound~$k$ is fixed ahead of time and must
be established before verification begins. This can be limiting: the cutoff
may be large, difficult to compute, or may not exist at all for systems with
richer interaction patterns or more complex topologies. Moreover, each cutoff
theorem covers a specific class of systems, and a protocol that falls slightly
outside the class, for instance because it combines conjunctive guards with
shared variables, is not covered.

\subsection{Dynamic Cutoffs and View Abstraction}

\emph{Dynamic cutoff methods} address these limitations by allowing the
abstraction boundary to be determined adaptively, during the verification
process itself, rather than fixed in advance. A well-developed technique in
this family is \textbf{view abstraction}~\cite{AHH13,AHH16}.
Rather than committing to a single small representative system upfront, the
method reasons about what a process can \emph{observe} about the rest of the
system, and uses those local observations to construct an on-the-fly
abstraction.

The starting point is to model the global state of an $n$-process system as a
configuration: a word of local process states, as in
Section~\ref{sec:prelim}. A \emph{view} of size~$k$ is a subword of length~$k$
of a configuration, that is, a projection of the configuration onto a subset
of~$k$ processes that preserves their order. Instead of tracking full
configurations (whose length varies with~$n$), the method tracks the set of
all views of size up to~$k$ that are reachable across \emph{any} instance of
the system.

For a given precision parameter $k \in \mathbb{N}$, the framework runs two
procedures in parallel. The first is an \emph{under-approximation}: it
performs concrete reachability analysis on instances of fixed size up to~$k$,
and reports a bug if a bad configuration is found. The second is an
\emph{over-approximation} based on abstract interpretation. We write
$\alpha_k$ for the abstraction function, which maps a configuration to its set
of subwords of size up to~$k$. We write $\gamma_k$ for the corresponding
concretization function, which maps a set of views to the set of all
configurations whose subwords of size up to~$k$ all belong to the given set of
views. A fixpoint $V_k$ is then computed over
views using an abstract post operator. The key soundness property is that the
set of reachable configurations~$R$ satisfies $R \subseteq \gamma_k(V_k)$ for
all~$k$, and the approximation becomes tighter as $k$ grows:
$\gamma_{k+1}(V_{k+1}) \subseteq \gamma_k(V_k)$.

An important observation for the implementation is that the abstract post
operator can be computed using only concrete configurations of size up to $k{+}1$, rather than
the full (infinite) concretization. This is because the witnesses required to
enable a global transition consist of a single extra process; accommodating
one additional process beyond the view size is always sufficient. This makes
the fixpoint computation entirely finite and automatic.

If the over-approximation establishes $\gamma_k(V_k) \cap \Bad = \emptyset$,
the system is declared safe. If the under-approximation finds a reachable bad
configuration, the system is declared unsafe. In either case, the current~$k$
is a \emph{dynamic cutoff}, a point discovered during verification beyond
which the search need not continue. If neither procedure is conclusive, the
precision is increased by incrementing~$k$ and the analysis is repeated.

\paragraph{The running example.}
Consider the running example with $k = 1$. The views of size one are single
local states, and all three local states $\Idle$, $\Wait$, and~$\Crit$ are
reachable. The concretization of this set of views contains every
configuration, including $\Crit\Crit$, so the over-approximation cannot prove
mutual exclusion, and the under-approximation, which only explores the
instance of size one, finds no bug. The precision is therefore increased to $k
= 2$.

With $k = 2$, the analysis starts from the views of the initial
configurations, namely $\Idle$ and $\Idle\Idle$. The abstract post operator
considers the configurations of size at most three whose views are all in the
current set, applies the transitions, and collects the views of the results.
In the first rounds, the transition from $\Idle$ to $\Wait$ produces the views
$\Wait$, $\Idle\Wait$, $\Wait\Idle$, and $\Wait\Wait$, and the transition from
$\Wait$ to $\Crit$ then produces $\Crit$, $\Idle\Crit$, $\Crit\Idle$,
$\Wait\Crit$, and $\Crit\Wait$. The view $\Crit\Crit$ is never produced: to
produce it, a process in~$\Wait$ would have to move to~$\Crit$ in a
configuration of size at most three that already contains a process
in~$\Crit$, and the guard forbids this. The fixpoint~$V_2$ therefore consists
of all words of length at most two except $\Crit\Crit$. Its concretization is
the set of configurations with at most one process in~$\Crit$, which is
disjoint from the bad configurations. The method thus proves mutual exclusion
for all sizes, with dynamic cutoff~$2$.

\paragraph{Termination and applications.}
The procedure is guaranteed to terminate whenever the set of reachable
configurations is downward-closed with respect to the subword ordering, as
defined in Section~\ref{sec:wqo}. It also terminates for \emph{monotonic
systems} in the sense of Section~\ref{sec:wqo} when the bad configurations
form an upward-closed set. In that case, safety depends only on the downward
closure of the reachable set, which is closed under the transitions and is
represented exactly by views of bounded size. Monotonic systems cover a broad
and well-studied class, including lossy channel systems, Petri nets (where the
method solves the coverability problem), and timed Petri nets. The running
example is not monotonic, but its reachable set, the set of configurations
with at most one process in~$\Crit$, is downward-closed, so the procedure
terminates as well. For Burns' mutual exclusion protocol, the method likewise terminates at~$k =
2$: already with two-process views it can be shown that no bad configuration
is reachable for any system size.

View abstraction has been successfully applied to a range of protocols,
including mutual exclusion algorithms such as the fine-grained version of
Szymanski's protocol, a benchmark that, to the knowledge of the authors
of~\cite{AHH16}, had not been verified automatically by other methods, as well as cache coherence protocols and distributed
systems with complex process interactions. The framework extends naturally to
other process topologies such as multisets, rings, and trees, and to settings
where global transitions are not assumed to be atomic~\cite{AHH13,AHH16}.
Further refinements of the abstraction, which take into account the context in
which a process takes a step, have been developed in~\cite{AHH14}.

Other dynamic approaches pursue the same goal by different means. Kaiser,
Kroening, and Wahl~\cite{KKW10} consider programs in which an unbounded number
of threads execute the same code over shared variables. They explore instances
of increasing size and detect, during the exploration, a point beyond which
additional threads cannot reach new local states. Like view abstraction, the
method determines the cutoff from the program at hand instead of from a
general theorem.

Together, static cutoff theorems and dynamic approaches like view abstraction
form a complementary toolkit: static cutoffs offer clean theoretical
guarantees for specific protocol families, while dynamic methods
provide greater flexibility and automation for protocols that fall outside
those families. The idea of reasoning about a parameterized system through
its small instances also reappears in the invisible invariants method of
Section~\ref{sec:deductive} and in the small-model properties of the decidable
logics used by SMT-based tools in Section~\ref{sec:smt}.
